
\documentclass[11pt]{article}
\usepackage[utf8]{inputenc}
\usepackage[spanish]{babel}
\usepackage{amsmath}
\usepackage{geometry}
\geometry{margin=1in}

\title{Documentación: Random Walk}
\date{}

\begin{document}
\maketitle

\section{Objetivo}
Estudiar el comportamiento de un proceso Random Walk en series de tiempo.

\section{Definición}
Un Random Walk puede expresarse como:
\[
X_t = X_{t-1} + \varepsilon_t
\]

donde $\varepsilon_t$ representa ruido blanco.

\section{Propiedades}
\begin{itemize}
\item No estacionario.
\item La varianza crece con el tiempo.
\item Presenta alta persistencia.
\item Los shocks tienen efectos permanentes.
\end{itemize}

\section{Importancia}
Muchos precios financieros pueden aproximarse mediante procesos Random Walk.

\end{document}
